% =====================================================================
%  CAPÍTULO II / 3 Metodologia
% =====================================================================
\setdiv{CAPÍTULO II}
\chapter{Metodologia}

A metodologia deste trabalho foi estruturada em etapas sequenciais, visando o
desenvolvimento, implementação e validação da ferramenta algorítmica proposta para a
negociação de ativos financeiros. O objetivo foi descrever os passos para construir e
testar a solução, desde a escolha do ambiente até a avaliação de seu desempenho.

\section{Natureza da Pesquisa}

Este trabalho, como mencionado na introdução, adota uma natureza de pesquisa
experimental. A pesquisa experimental, no contexto de desenvolvimento de sistemas e
algoritmos, envolve a manipulação controlada de variáveis que, neste caso, foram os
parâmetros dos modelos, as combinações de indicadores técnicos, para observar e analisar
os efeitos resultantes em um ambiente controlado. Foram conduzidos experimentos para testar
diferentes configurações da ferramenta algorítmica, buscando identificar as abordagens
mais eficazes para a negociação do ativo financeiro selecionado \cite{wazlawick2014}.

\section{Padrões para Pesquisa Experimental}

Para a condução dos experimentos e o desenvolvimento da ferramenta algorítmica
proposta, foi utilizado um conjunto definido de padrões e ferramentas, cujas
funcionalidades e conceitos foram detalhados no capítulo anterior de Fundamentação
Teórica.

A plataforma MetaTrader 5 (MT5) foi empregada como o ambiente central para as fases
de desenvolvimento e teste. A linguagem de programação MQL5, nativa do MT5, foi
utilizada para a codificação dos \textit{Expert Advisors} (\textit{EAs}). Estes EAs são os algoritmos
que efetivamente incorporaram a lógica dos modelos de análise de tendência,
volatilidade e modelagem condicional para a execução automatizada de ordens
de negociação.

O testador de estratégias integrado ao MT5 foi uma ferramenta crucial, permitindo a
realização de backtests. Esses testes simularam o desempenho operacional das
estratégias desenvolvidas ao longo de períodos históricos, utilizando os dados de
cotações passadas para avaliar a viabilidade e a rentabilidade potencial.

Como fonte primária de dados, foi utilizado o banco de dados histórico de cotações
fornecido pela própria plataforma MetaTrader 5 com integração a corretoras
brasileiras. Inicialmente, a conta era integrada à corretora Rico, que permitia acessar
ações da B3. No entanto, durante o desenvolvimento, essa corretora cessou o fornecimento
de ativos para a plataforma MetaTrader 5, necessitando uma migração para a corretora
Clear, que mantém compatibilidade com a plataforma. Esta escolha metodológica simplificou
significativamente o processo de aquisição de dados, seu tratamento inicial e a
subsequente manipulação, uma vez que os dados se encontram em um formato padronizado
e são diretamente acessíveis via MQL5. Os dados extraídos e utilizados incluíram os
preços de abertura, máximo, mínimo e fechamento, bem como o volume de negociação para
cada período do timeframe definido para a análise, que foi diário (D1) para capturar
tendências mais longas.

O ativo referente às ações preferenciais da Petrobras (PETR4) foi o objeto de estudo
principal. Sua notória volatilidade histórica é um aspecto central para a aplicação
dos modelos subsequentes, em particular o modelo GARCH, que é projetado
especificamente para capturar e modelar a variância condicional em séries temporais
que exibem tais flutuações. A liquidez e a representatividade do PETR4 no Ibovespa
foram consideradas para assegurar a disponibilidade de dados robustos para esta
análise de volatilidade e para contextualizar os resultados no mercado acionário
nacional.

Para uma caracterização inicial e visualização dessa volatilidade inerente ao PETR4,
foram calculadas e analisadas as Bandas de Bollinger. Este indicador técnico, que
constrói um canal em torno de uma média móvel utilizando o desvio padrão dos preços,
foi aplicado à série histórica de preços do ativo. A observação da largura das bandas e
do comportamento do preço em relação a elas forneceu insights valiosos sobre os
diferentes regimes de variação dos preços (períodos de alta ou baixa volatilidade) e
contribuiu para a compreensão do ambiente de mercado no qual a ferramenta algorítmica
operaria.

\section{Obstáculos Encontrados}

Durante o desenvolvimento deste trabalho, foram enfrentados obstáculos significativos que impactaram o escopo e a execução da pesquisa.

\textbf{Perda do ambiente Windows e limitações no Linux:} O ambiente primário de desenvolvimento era uma máquina Windows, equipada com o MetaTrader 5 instalado nativamente. A perda desse equipamento obrigou a migração para um ambiente Linux (EndeavourOS), no qual o MetaTrader 5 não possui suporte oficial. A solução adotada foi executar a plataforma via Wine, o que, embora funcional, impõe limitações operacionais, como instabilidades ocasionais e restrições de compatibilidade com algumas APIs nativas da plataforma.

\textbf{Incompatibilidade da corretora Rico:} Durante a fase de coleta de dados e execução de backtests com ativos de ações, a corretora Rico cessou a disponibilização de ações em sua plataforma MetaTrader 5, inviabilizando o acesso ao ativo PETR4 através dessa conta. Foi necessário abrir uma nova conta na corretora Clear, que mantém compatibilidade com a plataforma, para dar continuidade à pesquisa. Esse processo causou atrasos e exigiu a reconfiguração do ambiente de testes.

Esses fatores combinados contribuíram para a redução do eescopo original (que previa a integração de modelos de \textit{Machine Learning}) concentrando o trabalho nos modelos ARIMA e GARCH.

\section{Experimento de Pesquisa}

O experimento de pesquisa iniciou-se com a implementação de abordagens para análise
de tendência utilizando Médias Móveis. Foram programados em MQL5 diferentes tipos de
médias móveis, como a Simples e a Exponencial. Estratégias baseadas em cruzamentos
dessas médias foram implementadas para gerar indicações preliminares sobre a direção
potencial do mercado, ou seja, sinais de compra ou venda. Estas estratégias também
serviram para verificar o desempenho de modelos baseados apenas em médias no mercado
mais recente.

Posteriormente, o experimento avançou para uma modelagem estatística mais sofisticada
da série temporal de preços, utilizando o modelo ARIMA. A determinação da ordem de
integração `d' do modelo ARIMA $(p, d, q)$ foi realizada através da aplicação
sequencial do teste de raiz unitária ADF sobre a série. Para a definição das ordens
autorregressiva `p' e de médias móveis `q', foi adotado um procedimento sistemático
de testar múltiplas combinações de 0 a 4. A seleção da especificação final foi
baseada no critério de informação \textit{Akaike Information Criterion} (\textit{AIC}).

Na sequência, os resíduos obtidos a partir do modelo ARIMA ajustado foram utilizados
como série de entrada para a modelagem da volatilidade condicional, por meio de um
modelo GARCH. Esta abordagem combinada, ARIMA-GARCH, visou capturar a dinâmica da
variância dos erros que não foi explicada pelo modelo de média condicional ARIMA.
Foram estimados modelos GARCH, com foco inicial e principal na especificação GARCH
(1,1), dada a sua prevalência e bom desempenho demonstrado na literatura para séries
temporais financeiras. O objetivo foi modelar e prever a volatilidade futura de curto
prazo dos resíduos, fornecendo uma medida do risco associado às previsões de preço do
ARIMA.

\section{Diagramas}

\begin{figure}[htb]
\centering
\caption{Modelo Completo}
\label{fig:modelo-completo}
\includegraphics[width=0.45\textwidth]{fig08-modelo-completo.png}
\fonte{Autoria própria.}
\end{figure}

\section{Critérios para Avaliação}

Os testes das estratégias baseadas em Médias Móveis, cruzamento de Médias Móveis, e
dos modelos ARIMA e GARCH foram realizados integralmente na plataforma MetaTrader 5
por meio do ambiente de backtesting.

O desempenho dessas estratégias foi quantitativamente avaliado por meio de métricas
como a taxa de acerto das operações, o lucro ou prejuízo acumulado, o drawdown máximo
observado e o fator de lucro. Estas métricas foram fornecidas diretamente pela
plataforma. Adicionalmente, um critério importante para a validação interna dos modelos
estatísticos foi a qualidade do ajuste dos modelos ARIMA e GARCH. Esta qualidade foi
aferida utilizando o critério de informação Akaike (AIC) e pela verificação da
estacionariedade dos resíduos através do teste de raiz unitária de Dickey-Fuller
Aumentado (ADF).

\section{Resultados preliminares atingidos no PTG}

Como mencionado anteriormente, este projeto busca revisitar e expandir as investigações
de \citeonline{almeida2016}, que desenvolveu e testou um Expert Advisor (EA) baseado na
estratégia de cruzamento de Médias Móveis (MM) para o ativo PETR4 na plataforma
MetaTrader 5. O estudo original de \citeonline{almeida2016} observou que, embora sua
estratégia pudesse gerar lucros em determinados períodos de mercado mais estáveis, ela
enfrentava dificuldades e apresentava prejuízos durante crises financeiras,
evidenciando a sensibilidade do algoritmo simples às condições de mercado.

Para estabelecer um ponto de partida contemporâneo e avaliar o desempenho dessa
estratégia no cenário de mercado mais recente, foi necessário adaptar e reimplementar a
lógica central do EA proposto por \citeonline{almeida2016}.

O código original fornecido por ele, embora conceitualmente claro, demandou
atualizações para ser compatível com as versões mais recentes da linguagem MQL5 e da
plataforma MetaTrader 5. Adicionalmente, \citeonline{almeida2016} mencionou o uso de
bibliotecas externas para funcionalidades como gerenciamento de dinheiro e trailing
stops, as quais não foram disponibilizadas em seu trabalho. Desta forma, para a presente
análise preliminar, foi desenvolvido um EA em MQL5 utilizando a classe
CTrade\footnote{Biblioteca padrão do MQL5 para as operações de negociação.} padrão para
as operações de negociação e implementando a lógica de cruzamento de duas Médias Móveis
Simples (MMS).

A configuração do EA para este teste baseou-se nos seguintes parâmetros: Médias Móveis
Simples com período de 10 para a rápida e 20 para a lenta, aplicadas ao preço de
fechamento. A estratégia implementada inclui a reversão de posições, ou seja, ao ser
gerado um novo sinal, qualquer posição oposta existente é primeiramente fechada.

Os parâmetros de Stop Loss (SL) e Take Profit (TP) foram definidos por experimentação
inicial em 50 pontos e 100 pontos e depois nulos, respectivamente, uma vez que
\citeonline{almeida2016} não especificou os valores exatos utilizados em todos os seus
testes para PETR4. Reconhece-se que esta escolha de SL/TP pode divergir da utilizada
originalmente e constitui um ponto para otimizações e análises futuras.

As operações foram configuradas para serem executadas apenas na abertura de uma nova
barra no gráfico diário (D1), com a modelagem do backtest para \textit{Cada Tick} para maior
precisão na simulação dos preços de execução, alinhando-se com a abordagem de ``Every
tick'' mencionada por \citeonline{almeida2016}.

Visando uma comparação de desempenho sob condições financeiras similares às de um dos
cenários de \citeonline{almeida2016}, o backtest foi conduzido com um depósito inicial
de R\$ 50.000,00. O volume operacional foi definido como 1000 unidades do ativo PETR4
por negociação. A alavancagem utilizada na simulação foi de 1:100. O período de teste
compreendeu de 01 de janeiro de 2020 a 01 de maio de 2025.

Os principais resultados obtidos neste backtest para o ativo PETR4 são apresentados e
comparados qualitativamente com os achados de \citeonline{almeida2016} para PETR4 em
períodos de crise:

\begin{table}[htb]
\centering
\caption{Comparação modelo proposto x ALMEIDA (2016)}
\label{tab:comparacao-almeida}
\fontsize{10}{12}\selectfont
\begin{tabular}{lcccc}
\toprule
& \textbf{Pereira (2025)} & \textbf{Pereira (2025)} & \textbf{Almeida (2016)} & \textbf{Almeida (2016)} \\
& \textbf{1º teste} & \textbf{2º teste} & & \\
\midrule
\textbf{Período} & 2020-2025 (Crise) & 2020-2025 (Crise) & 2013-2016 & 2007-2011 (Crise) \\
\addlinespace
Volume por Operação & 1000 ações & 1000 ações & 1000 ações & 100 ações \\
Lucro Líquido Total & R\$ -9.230,00 & R\$ -2.530,00 & R\$ +10.430,00 & R\$ -2.542,00 \\
Retorno Percentual & -18,46\% & -5,06\% & +20,86\% & -5,08\% \\
Fator de Lucro & 0,64 & 0,93 & \textgreater{} 1,0 (lucro) & \textless{} 1,0 (prejuízo) \\
Total de Negociações & 64 & 58 & 38 & 52 \\
\bottomrule
\end{tabular}

\vspace{2pt}
\fonte{Elaboração própria com dados do MetaTrader 5 e pesquisa de \citeonline{almeida2016}.}
\end{table}

A análise destes resultados preliminares indica que, para o primeiro teste no período
recente de janeiro de 2020 a maio de 2025, a estratégia base de cruzamento de Médias
Móveis, operando com um volume de 1000 ações de PETR4, resultou em um prejuízo líquido
substancial de R\$ 9.230,00, representando uma perda de 18,46\% sobre o capital inicial.
O Fator de Lucro de 0,64 e a baixa taxa de acerto (apenas 25\% das 64 operações foram
lucrativas) confirmam a ineficácia da estratégia neste intervalo de tempo.

Em uma segunda rodada de testes preliminares para o período de janeiro de 2020 ao começo
de 2025 e mantendo o volume de 1000 ações de PETR4, foram zerados os parâmetros de Stop
Loss e Take Profit, resultando em um desempenho ligeiramente menos desfavorável. Neste
segundo cenário, o prejuízo líquido total foi de R\$ 2.530,00, correspondendo a uma
perda de 5,06\% sobre o capital inicial. Embora o Fator de Lucro tenha melhorado para
0,93 e o número total de negociações tenha sido reduzido para 58, a estratégia ainda não
demonstrou capacidade de gerar lucro. Este resultado, apesar de melhor que o anterior,
ainda se alinha com as dificuldades encontradas por \citeonline{almeida2016} em períodos
de crise e reforça a necessidade de abordagens mais robustas.

\begin{figure}[htb]
\centering
\caption{Desempenho do modelo proposto (2º teste)}
\label{fig:desempenho-2teste}
\includegraphics[width=0.8\textwidth]{fig09-desempenho-2teste.png}
\fonte{Elaboração própria com dados do MetaTrader 5.}
\end{figure}

\begin{figure}[htb]
\centering
\caption{Gráfico de dispersão da correlação de lucro (2º teste)}
\label{fig:dispersao-lucro}
\includegraphics[width=0.9\textwidth]{fig10-dispersao-lucro.png}
\fonte{Elaboração própria com dados do MetaTrader 5.}
\end{figure}

Este desempenho contrasta significativamente com os resultados obtidos por
\citeonline{almeida2016} para o mesmo ativo e volume em um período anterior (2013-2016),
onde foi reportado um lucro considerável. Por outro lado, o resultado negativo atual se
assemelha, em natureza, aos prejuízos que Almeida também encontrou ao aplicar sua
estratégia em PETR4 durante períodos de crise financeira, como a de 2008.

O período de 2020-2025 foi caracterizado por uma volatilidade acentuada e por dinâmicas
de mercado complexas, influenciadas por eventos globais como a pandemia de COVID-19 e
suas repercussões econômicas. É plausível que tais condições tenham sido
particularmente desfavoráveis para uma estratégia de seguimento de tendência simples
como o cruzamento de médias. Adicionalmente, as necessárias adaptações do código
original de \citeonline{almeida2016} e a definição experimental dos parâmetros de Stop
Loss e Take Profit, na ausência dos valores exatos utilizados no estudo de referência,
são fatores que podem ter influenciado os resultados.

No entanto, a principal inferência destes testes preliminares é a clara demonstração de
que a estratégia de cruzamento de Médias Móveis, mesmo com parâmetros de risco básicos
e um volume operacional mais elevado, não se mostrou suficiente para gerar retornos
positivos de forma consistente no mercado recente de PETR4. Este achado é fundamental,
pois estabelece um baseline contemporâneo e valida de forma contundente a justificativa
central deste trabalho de graduação: a necessidade premente de investigar, desenvolver
e integrar modelos analíticos e preditivos mais sofisticados, como ARIMA para previsão
de tendência e GARCH para modelagem da volatilidade, com o objetivo de construir uma
ferramenta algorítmica mais robusta, adaptável e com maior potencial de eficácia para a
negociação de ativos voláteis no dinâmico cenário atual do mercado de ações brasileiro.

\section{Limitações e Trabalhos Futuros}

Os resultados finais de backtesting com os modelos ARIMA e GARCH ainda estão sendo consolidados e analisados. Os resultados preliminares obtidos com estratégias de Médias Móveis demonstraram a ineficácia de abordagens simplistas em cenários de mercado volátil, validando a necessidade de modelos mais sofisticados.

As limitações encontradas durante este trabalho incluem a restrição computacional e operacional do ambiente Linux via Wine, bem como as interrupções causadas pela mudança de corretora. Adicionalmente, o escopo reduzido — concentrado em ARIMA e GARCH — significa que componentes originalmente planejados, como a integração de Machine Learning (LSTM) para análises não lineares, não foram implementados neste ciclo.

Como perspectivas para trabalhos futuros, recomenda-se:

\begin{itemize}
  \item Implementação do modelo LSTM em um ambiente de desenvolvimento dedicado (Python com TensorFlow/Keras), seguido de integração ao MetaTrader 5 como componente complementar aos modelos ARIMA e GARCH;
  \item Incorporação de indicadores técnicos como filtro de confirmação (RSI, MACD) para validar sinais gerados pelos modelos estatísticos;
  \item Expansão da análise para outros ativos brasileiros de alta liquidez (ex., BBDC4, PETR3, VALE3) para validar a generalização da metodologia;
  \item Otimização de hiperparâmetros dos modelos ARIMA e GARCH através de algoritmos evolucionários (Genetic Algorithm, PSO);
  \item Implementação de técnicas robustas de gerenciamento de risco (trailing stops, position sizing dinâmico) para mitigação de drawdowns.
\end{itemize}
